% !TeX root = ../../金融经济学_王江_讲义及习题.tex
% 知识点：references/金融经济学/金融经济学-王江_261013.pdf，印刷页 1–5（PDF 第 16–20 页）
% 补充：references/金融经济学/金融经济学_第1章 引论.pptx，第 7–46 页
% 本章无原书习题；PPT 的教学组织、考核安排与联系方式不纳入正文。

\lecture{引论}

本章先说明金融经济学研究什么，再区分均衡定价与无套利定价，最后用一个两状态期权例子展示如何从明确假设推导可检验的价格约束。教材的引论给出全书路线；PPT 进一步讨论学科发展、数学公理化方法与模型的适用边界。

\begin{tcolorbox}[knowledge,title=1.1 金融经济学的研究对象]
\textbf{金融学研究不确定环境下资源的跨期配置。}今天与未来的资源不同，未来不同状态下的资源也不同；金融交易使参与者能够交换对这些资源的要求权。金融市场既是金融要求权（financial claim）的交易场所，也是形成金融资产价格、引导资源配置的机制。

资源配置有两个相关层面：\textbf{微观层面}，个人、公司等如何利用既有资源满足需求；\textbf{整体层面}，不同参与者之间如何转移资源与分担风险。个人在既定约束下最优，并不自动意味着全社会配置有效，后者还取决于市场结构与交易条件。

\textbf{教材的三个核心问题：}
\begin{enumerate}
 \item 参与者如何作出金融决策，尤其是交易与组合选择？
 \item 这些选择如何共同决定市场行为与金融要求权的价格？
 \item 价格如何影响资源的实际配置，市场能否实现有效配置？
\end{enumerate}
金融经济学用经济学的一般原理分析这些问题，为投资、金融市场与机构、金融工程以及公司财务提供理论基础。金融资产的突出特点是\textbf{面向未来的支付要求权}；其价格取决于未来支付、信息、风险及参与者的风险态度。金融学与经济学并非两套互斥的方法，二者共享选择与均衡分析，金融学尤其重视跨期、不确定性与无套利约束。
\end{tcolorbox}

\begin{tcolorbox}[knowledge,title=1.2 均衡定价与无套利定价]
\textbf{均衡定价：}从经济的基本要素出发，包括资源禀赋、偏好、未来风险和交易机会，在价格给定时求每个人的最优选择，再寻找使市场出清的价格。\textbf{最优化与出清须同时成立}；仅仅所有人各自没有意愿继续调整，并不足以代替出清条件。

\textbf{无套利定价：}给定若干已有证券的价格与支付，约束其他相关证券的价格。如果一个支付能由已有证券组合逐状态复制，则其价格必须等于复制成本；否则买入便宜的一方、卖出昂贵的一方即可套利。\textbf{复制是相对定价的核心技术。}

\noindent\begin{tabularx}{\linewidth}{@{}lXX@{}}
\toprule
 & 均衡定价 & 无套利定价\\
\midrule
给定条件 & 偏好、禀赋、风险、市场结构 & 已有证券价格、支付与可行交易\\
基本要求 & 每人最优且市场出清 & 不存在无风险的免费净收益\\
回答问题 & 基础价格与配置如何共同形成 & 相关支付的价格如何保持一致\\
局限 & 结论依赖经济假设与均衡条件 & 不完全市场中未必确定唯一价格\\
\bottomrule
\end{tabularx}

PPT 将二者分别称作“绝对定价”与“相对定价”。这里“绝对”不是脱离计价单位的数值：均衡价格仍须选定计价物。无套利也不意味着偏好完全不影响价格，而是\textbf{偏好对已有价格的影响已被包含在这些价格中}，复制定价时无须再次估计。

在本书的无摩擦与适当单调性条件下，均衡必无套利；反过来，无套利通常只给出价格一致性的必要约束，并不保证市场已经出清。第四章给出正式定义与基本定理，第五章展开期权定价。
\end{tcolorbox}

\begin{tcolorbox}[knowledge,title=1.3 基本框架与全书路线]
教材采用一个简化的\textbf{静态、两时点证券市场框架}：以个人参与者为原型，研究外生给定支付的证券，先把市场结构视为给定。静态并不是没有未来或没有风险，而是不展开中间日期不断接收信息、调仓交易的过程。

\textbf{暂时抽掉的两个关键因素是信息不对称与交易成本。}现实金融要求权可能依赖行为、难以完全写入合约，交易也可能通过中介而非公开证券市场进行；机构与市场结构还会随技术、激励和制度变化。本书先建立基准，不表示这些因素不重要，更不表示基准结论在存在摩擦时可以直接照搬。

\noindent\begin{tabularx}{\linewidth}{@{}lX@{}}
\toprule
章节 & 依次解决的问题\\
\midrule
第 2–3 章 & 建立经济环境、偏好与市场模型；在 Arrow--Debreu 市场分析选择、均衡与配置效率。\\
第 4–5 章 & 从一般市场结构推出无套利约束，再用于期权定价与市场完全化。\\
第 6–9 章 & 期望效用、风险厌恶与投资组合选择，进一步分析参与者行为。\\
第 10–11 章 & 完全和不完全市场中的均衡及资产定价。\\
第 12–14 章 & 以均值—方差偏好及因子结构获得 CAPM、APT 等更具体模型。\\
第 15 章 & 将框架用于公司财务，包括价值最大化、所有权与管理权分离、MM 定理。\\
\bottomrule
\end{tabularx}

教材结尾再讨论信息不对称、交易成本等被简化因素，附录补充证明和数学工具。\textbf{建模的顺序是先给出一般基准，再增加具体结构}，而不是把每章的公式看成相互独立的经验规则。
\end{tcolorbox}

\begin{tcolorbox}[knowledge,title=1.4 学科发展的主要线索（PPT 补充）]
PPT 用若干历史节点说明工具如何形成：Bernoulli 的不确定性下效用分析（1738）、Bachelier 的价格随机运动研究（1900）、Fisher 的跨期选择与分离思想（1930）、Keynes 的流动性偏好（1936），以及 von Neumann--Morgenstern 的期望效用公理体系（1944）。这些分别回应了\textbf{如何比较风险选择、如何描述价格变化、如何处理跨期配置、如何分析货币利率和如何建立一致的决策标准}。

此后，投资组合理论、一般均衡理论、无套利与期权定价逐步形成现代微观金融学的分析框架。Arrow--Debreu 的一般均衡研究与 Debreu 的《价值理论》强调用集合、偏好和凸性来明确可行选择与价格体系。组合选择、CAPM、期权定价、MM 定理与 APT 则把这些思想用于不同的金融问题。

PPT 概括的三个变化是：分析工具更加定量化，分析基础更加重视个体决策，理论逐步成为设计与评估金融安排的工具。\textbf{历史联系不应替代具体假设：}例如鞅说的是条件期望关系，并不推出“最可能的未来价格等于现在价格”；后者混淆了均值与众数，也忽略了计价单位和所用概率测度。
\end{tcolorbox}

\begin{tcolorbox}[knowledge,title=1.5 数学公理化：假设、推导与经验检验]
公理化方法先定义对象、可行选择和行为假设，再按逻辑推导结论。其作用是把“为什么成立”与“在哪些条件下成立”写清楚，使增加、放松某项假设时，能够识别结论随之发生的变化。

PPT 围绕几种方法论立场展开讨论，以下保留争论的要点，而不把评价性表述当作数学定理：
\begin{itemize}
 \item \textbf{Keynes 式的提醒：}符号复杂不等于解释深刻，推导仍依赖初始假设，不能因形式精确而忽略经济含义。
 \item \textbf{Debreu 强调的严格性：}明确对象和推理，能发现不一致、复用已证结果，并为研究新问题提供基础。
 \item \textbf{Popper 强调的可检验与可证伪：}理论应推出可能受经验反驳的结论；从假设演绎出来不等于已经得到数据支持。
 \item \textbf{Friedman 强调的预测表现：}评价理论要重视其解释与预测能力；但这不能被理解为所有不现实的假设都无须关注。假设决定模型可能在哪些环境下失效。
\end{itemize}
演绎推理回答“若假设成立，结论是否成立”；统计分析、计量和案例证据帮助回答“这些机制在观察对象中是否起作用”。样本中观察到的规律仍有外推与识别问题，不能被当作无条件的逻辑必然。模型、数据图表与案例可互相补充，不能用一种证据完全替代另一种。

\textbf{CAPM 的假设清单示例。}PPT 列出相同单期投资期限、资产可分且可交易、允许卖空、均值—方差选择、对收益均值与协方差的同质预期、同一利率下自由借贷、无税费与价格接受者等条件。列出它们的目的，是把定价结论与其成立环境对应起来；现实应用须检查约束、估计误差和模型偏离，而非只记一条收益公式。

\textbf{实证与规范。}实证分析研究“是什么、为何如此”，规范分析还须引入“应当如何”的价值标准。经济学研究和投资实践都可能使用这两类分析，不能按职业作绝对划分。配置有效与否也不同于分配公平与否，第二章将给出 Pareto 标准。
\end{tcolorbox}

\begin{tcolorbox}[example,title=PPT 推理例子：公共信息与共同知识]
将 PPT 的夫妻推理题写成等价的红标问题：$50$ 人头上各有一个红标，每人看得见别人的标记、看不见自己的。所有人都知道彼此严格按逻辑推理；每天统一报告，能确定自己有红标的人当天报告，所有人都能观察前一天的报告。起初没有公共宣告，后来一位可靠的主持人\textbf{公开宣布“至少一人有红标”}，所有人都知道所有人听到了这一宣告。

\textbf{结论：}从宣告后的第一天计，若实际有 $n$ 个红标，这 $n$ 人将在第 $n$ 天一起报告；因此本题为第 $50$ 天。

\textbf{归纳推导：}只有一个红标时，该人看不到其他红标，由宣告立即确定自己有标，在第一天报告。若有两个，每人看到一个；第一天无人报告，就排除了“只有对方一个”的可能，第二天确定自己也有标。一般地，假设 $n-1$ 个红标会在第 $n-1$ 天报告，实际有 $n$ 个时，每名红标者看到 $n-1$ 个；若自己无标，他们应按上述规则报告。观察到第 $n-1$ 天仍无人报告后，才能排除该可能，在第 $n$ 天报告。

当实际人数很多时，每个人在宣告前早已看到至少一个红标；宣告新增的是\textbf{所有人知道、知道别人知道、并如此递推的共同知识}，从而启动上述归纳。相同的一阶事实不等于相同的高阶知识结构；信息公开、观察规则和理性假设都是结论的一部分。
\end{tcolorbox}

\begin{tcolorbox}[knowledge,title=1.6 抽象时必须保留哪些差异]
PPT 用“外延与内涵”提醒我们：把证券仅描述为支付，会抽掉法律、流动性、交易约束等特征。应精确区分：\textbf{集合具有相同成员意味着集合相等，但并不意味着集合中任意两个成员具有相同性能、可以互换。}

无套利复制要求\textbf{同一日期、每个状态的净支付相同}，同时能按所用价格执行买卖；仅有相同的收益分布或期望收益还不够。例如两个等概率状态下，支付 $[1,0]^{\mathsf T}$ 与 $[0,1]^{\mathsf T}$ 的分布相同，却发生在不同状态，其价格可分别为 $\phi_1$ 与 $\phi_2$，不必相等。

期权与复制组合在特定模型下可以有相同支付，但现实中可能有不同的融资、保证金、税费或对手方安排。此时应扩充模型中的成本与约束，而非直接否定复制逻辑。反之，模型的逻辑自洽也不能替代适用条件的核对。
\end{tcolorbox}

\begin{tcolorbox}[knowledge,title=1.7 两状态期权：三种定价方法的同一逻辑]
沿用 PPT 的符号，股票现价为 $S_0>0$，未来为 $aS_0$ 或 $bS_0$，设 $a>b$。执行价为 $K$ 的欧式看涨期权支付为
\[
 C_a=[aS_0-K]_+,\qquad C_b=[bS_0-K]_+.
\]
其中 $[x]_+=\max\{x,0\}$。设一期无风险总收益率为 $R=1+r_F>0$，交易无摩擦、可卖空且两状态均可能发生。股票与债券无套利要求 $b<R<a$：否则股票与按无风险利率融资的组合会有零成本、非负非零的净支付。

\textbf{A. 复制组合。}持有 $z$ 股股票与一笔当前投入 $y$ 的无风险账户，未来支付为 $zaS_0+Ry$、$zbS_0+Ry$。逐状态匹配期权支付：
\[
 zaS_0+Ry=C_a,\qquad zbS_0+Ry=C_b.
\]
两式相减并代回，得到
\[
 z=\frac{C_a-C_b}{S_0(a-b)},\qquad
 y=\frac{aC_b-bC_a}{R(a-b)},\qquad C_0=zS_0+y.
\]
期权与复制组合未来相同，所以初始成本相同；若报价偏高就卖出期权、买入复制组合，偏低则反向操作。

\textbf{B. 风险对冲。}当 $C_a\ne C_b$ 时，卖空一股股票、买入 $x$ 份期权，取
\[
 x=\frac{S_0(a-b)}{C_a-C_b}=\frac1z,
\]
使两状态支付相同。无风险组合必须按 $R$ 增长，故
\[
 R(-S_0+xC_0)=-aS_0+xC_a=-bS_0+xC_b.
\]
解得与复制法相同的价格。若 $C_a=C_b$，不能用此处含零分母的对冲比，而应直接把期权视为确定支付，价格为 $C_a/R$。

\textbf{C. 风险中性定价。}定义 $q=(R-b)/(a-b)$，则 $0<q<1$，并有 $qa+(1-q)b=R$。将复制价格整理为
\[
 \boxed{C_0=\frac{qC_a+(1-q)C_b}{R}.}
\]
这不是假设投资者真的风险中性，也不是估计真实上涨概率，而是将已有价格约束改写为另一种期望。按无风险账户计价后，价格在该测度下满足鞅关系。

\textbf{与 PPT 的公式对应：}第 43–45 页公式没有写利息，等价于本例取 $R=1$，此时须有 $b<1<a$，$q=(1-b)/(a-b)$，且
\[
 C_0=\frac{C_a-C_b+aC_b-bC_a}{a-b}.
\]
若利率非零，就必须保留贴现因子；不能把未经规范化的价格直接说成等于未来价格的期望。
\end{tcolorbox}

\begin{tcolorbox}[knowledge,title=本章小结]
\begin{enumerate}
 \item 金融经济学研究不确定条件下的金融决策、价格形成与资源配置，核心是跨期配置和风险分担。
 \item 均衡定价从偏好、禀赋与市场结构推出最优化且出清的价格；无套利定价从已有价格约束相关支付的价格。
 \item 数学推导须写清假设，经济解释须与公式一致，经验应用还须检验模型与数据的对应关系。
 \item 共同知识不同于人人分别知道；逐状态相同支付也不同于相同概率分布。
 \item 两状态期权的复制、对冲与风险中性方法给出同一价格；PPT 的无贴现公式是零利率特例。
\end{enumerate}
\end{tcolorbox}
